% ============================================================
% Cheat sheet, argomento: i pattern a confronto e gli altri pattern della GoF (dispensa, capitolo 15)
% ============================================================
\argomento{Pattern a confronto}{Dispensa: capitolo 15. Quelli che si somigliano, e gli altri della GoF nominati
nelle slide}[Pattern a confronto]

\refstepcounter{voce}\label{v:confronti}
\begin{guida}{{\badge[pcolore]{\thevoce}\enspace I pattern che si somigliano}}
\renewcommand{\arraystretch}{1.3}
\begin{tabular}{@{}>{\bfseries\raggedright\arraybackslash}p{3.2cm} P{14.4cm}@{}}
  State o Strategy? & Stessa struttura (un contesto con un campo di tipo astratto e più implementazioni). Nello
    \textbf{State} sono gli \textbf{stati} a cambiare l'oggetto puntato (\code{setState}), si conoscono fra loro e
    rappresentano \textbf{fasi} della vita del contesto; nello \textbf{Strategy} lo sceglie il \textbf{client}, le
    strategie non si conoscono e sono \textbf{modi diversi} di fare la stessa cosa. \\
  Decorator o Composite? & Entrambi implementano il Component \textbf{e} ne contengono. Il \textbf{Decorator} ne
    contiene \textbf{uno} e \textbf{aggiunge} comportamento (una catena); il \textbf{Composite} ne contiene \textbf{molti}
    e \textbf{combina} i risultati dei figli (un albero). \\
  Decorator o Adapter? & Entrambi avvolgono un oggetto e gli delegano il lavoro. Il \textbf{Decorator} ha la
    \textbf{stessa} interfaccia dell'oggetto avvolto e aggiunge funzioni; l'\textbf{Adapter} avvolge un oggetto con
    un'interfaccia \textbf{diversa} e la traduce, senza aggiungere nulla. \\
  Decorator o Proxy? \oltre & Stessa forma. Il Decorator \textbf{aggiunge} responsabilità e si concatena; il Proxy
    \textbf{controlla l'accesso} all'oggetto vero (permessi, cache, creazione pigra) e di solito è uno solo. \\
  Adapter o Facade? \oltre & L'Adapter adatta \textbf{una} classe a un'interfaccia che esiste già; la Facade inventa
    un'interfaccia \textbf{semplice} davanti a un intero sottosistema di classi. \\
  Observer o Composite? & Entrambi hanno una lista e chiamano un metodo su tutti. Gli \textbf{osservatori} hanno
    un'interfaccia diversa da quella del soggetto e vengono avvisati di un cambiamento; i \textbf{figli} di un
    composto hanno la sua stessa interfaccia e formano un albero. \\
  Builder o costruttore? & Il costruttore riceve tutto in una volta (e con molti parametri facoltativi diventa
    illeggibile); i \code{set} lasciano l'oggetto a metà. Il \textbf{Builder} raccoglie i dati a passi e crea il
    prodotto, già controllato e immutabile, solo con \code{build()}. \\
  Singleton, classe statica o DI? & La classe con soli metodi \code{static} non ha un oggetto; il
    \textbf{Singleton} ha un oggetto unico che chi lo usa \textbf{va a prendere} (\code{getInstance()}); con la
    \textbf{Dependency Injection} l'oggetto (unico o no) viene \textbf{passato} a chi lo usa, e nei test si
    sostituisce. \\
  Strategy o DI? & Non si escludono: lo Strategy dice \textbf{che cosa} è l'oggetto (un algoritmo intercambiabile),
    la DI dice \textbf{come arriva} (da fuori, nel costruttore). Una strategia iniettata nel costruttore è entrambe
    le cose. \\
\end{tabular}
\end{guida}

\refstepcounter{voce}\label{v:altri-gof}
\begin{guida}{{\badge[pcolore]{\thevoce}\enspace Gli altri pattern della GoF, nominati nelle slide}}
Le slide li elencano con la loro definizione (seconda colonna); come riconoscerli nel codice (terza colonna) è
un'aggiunta \oltre, per non restare senza parole se un esercizio ne mostra uno.

\smallskip
\renewcommand{\arraystretch}{1.3}
\begin{tabular}{@{}>{\bfseries\raggedright\arraybackslash}p{2.6cm} P{6.6cm} P{8.3cm}@{}}
  Pattern & \textbf{Che cosa fa} & \textbf{Come si riconosce nel codice} \\ \hline
  Abstract Factory & Un'interfaccia per creare \textbf{famiglie} di oggetti correlati senza specificarne le classi
    concrete & un'interfaccia con più metodi \code{creaX()}, \code{creaY()}; una fabbrica concreta per famiglia
    (\code{TemaChiaro}, \code{TemaScuro}); il client riceve la fabbrica e non fa \code{new} \\
  Factory Method & Un'interfaccia per creare un oggetto, ma le \textbf{sottoclassi} scelgono la classe da istanziare &
    un metodo astratto \code{creaProdotto()} nella superclasse, ridefinito nelle sottoclassi; la superclasse usa il
    prodotto senza sapere quale sia \\
  Prototype & Crea i nuovi oggetti \textbf{copiando} un'istanza prototipo & \code{clone()} o \code{copia()} chiamato su
    un oggetto esistente al posto di \code{new}; spesso un registro di prototipi \\
  Bridge & \textbf{Disaccoppia} un'astrazione dalla sua implementazione, così che possano variare indipendentemente
    & due gerarchie parallele; le classi della prima hanno un campo del tipo della seconda (una \code{Forma} con un
    \code{Disegnatore}) \\
  Facade & Un'interfaccia \textbf{unificata} a un sottosistema, più facile da usare & una classe con pochi metodi
    semplici, ognuno dei quali chiama molte classi del sottosistema \\
  Proxy & Un \textbf{surrogato} di un oggetto, per controllarne l'accesso & stessa interfaccia dell'oggetto vero, un
    campo verso di lui, controlli prima di delegare (permessi, cache, caricamento pigro) \\
  Flyweight & La \textbf{condivisione} per gestire in modo efficiente molti oggetti piccoli & una fabbrica con una
    mappa che restituisce oggetti immutabili già creati invece di crearne di nuovi \\
  Template Method & Lo \textbf{scheletro} di un algoritmo, con alcuni passi definiti dalle sottoclassi & un metodo
    (spesso \code{final}) nella superclasse astratta che chiama, in ordine, metodi astratti o \code{protected}
    ridefiniti nelle sottoclassi \\
  Visitor & Un'operazione sugli elementi di una struttura; se ne aggiungono senza modificare le classi degli
    elementi & \code{accept(Visitor v)} negli elementi, che chiama \code{v.visit(this)}; un'interfaccia \code{Visitor}
    con un \code{visit} per ogni tipo di elemento \\
  Iterator & Accesso in sequenza agli elementi di una collezione (non è nelle slide) & \code{hasNext()} e
    \code{next()}, \code{Iterable}, il ciclo \code{for (x : collezione)} \\
  Command & Una richiesta trasformata in un oggetto (non è nelle slide) & un'interfaccia con \code{execute()}, oggetti
    comando messi in coda o in una lista, spesso con \code{undo()} \\
\end{tabular}
\end{guida}

\begin{attenzione}{I pattern non risolvono tutti i problemi}
Le slide chiudono con una provocazione: «i design pattern risolvono tutti i vostri problemi di ingegneria del
software». Non è così: prima si definisce il \textbf{problema}, poi si cerca il pattern che lo risolve. Un pattern
applicato dove non serve aggiunge solo classi e complessità; in un esercizio, il nome del pattern va sempre
\textbf{motivato} con il problema del testo e con la struttura del codice.
\end{attenzione}
